\chapter{Vitamin B6 (Pyridoxine)}
\section*{What Is This Test?}
Vitamin B6 refers to a group of three related pyridine derivatives: pyridoxine (PN), pyridoxal (PL), and pyridoxamine (PM), and their phosphorylated derivatives, primarily pyridoxal 5'-phosphate (PLP), the biologically active coenzyme form. PLP serves as a cofactor for over 140 enzymatic reactions, including amino acid metabolism (transamination, decarboxylation, deamination), neurotransmitter synthesis (serotonin, dopamine, norepinephrine, GABA), heme synthesis (delta-aminolevulinic acid synthase, a PLP-dependent enzyme), glycogenolysis (glycogen phosphorylase), and sphingolipid synthesis. Vitamin B6 deficiency classically causes seborrheic dermatitis, glossitis, cheilitis, microcytic hypochromic anemia (from impaired heme synthesis), peripheral neuropathy, confusion, depression, and seizures (infantile B6-dependent seizures). Deficiency is most commonly seen in alcoholism, malnutrition, older adults, and patients on isoniazid (which binds PLP and increases renal excretion), penicillamine, or oral contraceptives. The most clinically relevant biomarker is plasma pyridoxal 5'-phosphate (PLP).
\section*{Alternative Names}
Pyridoxal 5'-Phosphate; PLP; Pyridoxine; Vitamin B6 Level
\section*{Principle}
Plasma PLP is measured by HPLC with fluorescence detection, LC-MS/MS, or enzymatic assay using the PLP-dependent tyrosine decarboxylase apoenzyme. Specimen: plasma collected in an EDTA tube; the specimen should be protected from light because PLP is light-sensitive.
\section*{History}
Vitamin B6 was discovered in 1934 by Paul Gy\"orgy, who isolated a factor that cured acrodynia, a dermatitis, in rats; pyridoxine was crystallized and its structure identified by 1939. In 1945, Gunsalus and colleagues identified pyridoxal 5'-phosphate as the active coenzyme, explaining its central role in amino acid transamination. When isoniazid was introduced for tuberculosis in the 1950s, Biehl and Vilter showed in 1954 that it caused clinical pyridoxine deficiency with peripheral neuropathy, establishing routine vitamin B6 supplementation with isoniazid. In the same year, Hunt and colleagues described pyridoxine-dependent seizures in infants, an autosomal recessive disorder traced in 2006 to mutations in ALDH7A1 (antiquitin). In 1983, Schaumburg and colleagues described the irreversible sensory neuronopathy caused by megadose pyridoxine supplementation, defining the toxicity end of the spectrum.
\section*{Why This Test Is Ordered}
\begin{itemize}
  \item Evaluation of suspected vitamin B6 deficiency in patients with peripheral neuropathy, microcytic or sideroblastic anemia, seborrheic dermatitis, glossitis, or cheilitis
  \item Workup of unexplained hyperhomocysteinemia, since PLP is a cofactor for homocysteine catabolism
  \item Assessment of deficiency risk in alcoholism, malabsorption, older adults, and patients taking isoniazid or penicillamine
  \item Evaluation of infantile seizures when pyridoxine-dependent epilepsy (antiquitin deficiency) is suspected
  \item Monitoring of patients with chronic renal failure or on dialysis, who lose PLP with treatment
  \item Documentation of pyridoxine toxicity in patients with sensory neuropathy taking high-dose supplements ($>$200 mg/day)
\end{itemize}
\section*{Primary vs Secondary Test}
Plasma PLP measurement is the primary direct biomarker of vitamin B6 status. It is used secondarily to explain clinical syndromes (anemia, neuropathy, dermatitis) and to interpret functional markers such as elevated homocysteine.
\section*{How This Test Is Performed}
A blood sample is collected by standard venipuncture into an EDTA tube. The specimen is protected from light, because PLP is light-sensitive, and transported promptly to the laboratory. Plasma is separated and PLP is measured by HPLC with fluorescence detection, LC-MS/MS, or an enzymatic assay using the PLP-dependent tyrosine decarboxylase apoenzyme. Results are reported in nmol/L or ng/mL.
\section*{What Preparation Is Needed}
No fasting is required for plasma PLP measurement. Recent vitamin B6 supplementation should be disclosed because it raises PLP and can mask a true tissue deficiency; no supplement changes are needed when toxicity is suspected. The specimen must be protected from light and processed promptly.
\section*{Contraindications and Precautions}
\begin{itemize}
  \item No contraindications to standard venipuncture. Precautions:
  \item (1) The sample must be shielded from light and processed promptly because PLP degrades on exposure;
  \item (2) hemolyzed specimens are unreliable because red blood cells contain substantial PLP that falsely raises the measured level;
  \item (3) renal failure and hypoalbuminemia alter PLP levels and complicate interpretation; and
  \item (4) recent B6 supplementation elevates PLP and may mask deficiency.
\end{itemize}
\section*{Risks}
Minimal risk. Slight pain or bruising at the venipuncture site; rarely hematoma, infection, or vasovagal syncope.
\section*{What Happens After the Test}
Results are typically available within a few days. A low PLP (<20 nmol/L) is treated with oral pyridoxine replacement, and deficiency symptoms generally improve within weeks. Patients on isoniazid receive prophylactic pyridoxine 25--50 mg/day regardless of measured levels, and repeat testing after 1--3 months of therapy can document normalization.
\section*{Understanding Results}
Deficient: PLP <20 nmol/L (or <5 ng/mL). Causes: alcoholism, malnutrition, isoniazid therapy (isoniazid binds PLP and forms hydrazone complexes that are renally excreted), penicillamine, oral contraceptives, chronic renal failure (dialysis removes PLP), and inborn errors (pyridoxine-dependent epilepsy due to ALDH7A1/antiquitin deficiency). Toxic: excessive B6 supplementation (>200 mg/day for months) causes irreversible sensory neuropathy (pure sensory neuronopathy/ganglionopathy with ataxia and loss of proprioception).
\section*{Normal Results}
Plasma PLP: 20--175 nmol/L (5--40 ng/mL). Optimal: >30 nmol/L. Deficiency: <20 nmol/L. Elevated: >175 nmol/L (supplementation).
\section*{Follow-up Tests}
Serum homocysteine (B6 deficiency elevates), CBC (microcytic anemia), serum transaminases (AST/ALT are PLP-dependent, levels may be low in B6 deficiency), tryptophan loading test (elevated xanthurenic acid in urine after tryptophan load). Isoniazid-related deficiency: prophylactic pyridoxine 25--50 mg/day.
\section*{References}
\begin{enumerate}
  \item Leklem JE. Vitamin B-6: a status report. J Nutr. 1990;120(Suppl 11):1503--1507.
  \item Schaumburg H, Kaplan J, Windebank A, et al. Sensory neuropathy from pyridoxine abuse. N Engl J Med. 1983;309(8):445--448.
  \item Ueland PM, Ulvik A, Rios-Avila L, et al. Direct and functional biomarkers of vitamin B6 status. Annu Rev Nutr. 2015;35:33--70.
  \item Institute of Medicine. Dietary Reference Intakes for Thiamin, Riboflavin, Niacin, Vitamin B6, Folate, Vitamin B12, Pantothenic Acid, Biotin, and Choline. Washington, DC: National Academies Press; 1998.
  \item Clayton PT. B6-responsive disorders: a model of vitamin dependency. J Inherit Metab Dis. 2006;29(2--3):317--326.
\end{enumerate}
\clearpage
\section*{Regulatory Status and Authorization Date}
This chapter describes a test category, not a specific commercial product. FDA, CDSCO, EU IVDR, and MHRA approval or registration dates therefore cannot be assigned without the assay manufacturer, product name, instrument, and intended use. For laboratory-developed tests, an individual agency approval date may not exist. See the Preface for the interpretation of regulatory status.

\clearpage
